\documentclass[11pt,letterpaper]{article}

% -------- Packages --------
\usepackage[margin=1in]{geometry}
\usepackage{graphicx}
\usepackage{xcolor}
\usepackage{titlesec}
\usepackage{enumitem}
\usepackage{tabularx}
\usepackage{array}
\usepackage{booktabs}
\usepackage{longtable}
\usepackage{fancyhdr}
\usepackage{amsfonts}
\usepackage{hyperref}
\usepackage{parskip}
\usepackage{tcolorbox}
\tcbuselibrary{skins,breakable}

% -------- Colours --------
\definecolor{uavblue}{RGB}{25,72,120}
\definecolor{uavgrey}{RGB}{70,70,70}
\definecolor{uavlight}{RGB}{240,244,250}
\definecolor{uavwarn}{RGB}{170,60,30}
\definecolor{uavgood}{RGB}{40,110,60}

% -------- Hyperref --------
\hypersetup{
  colorlinks=true,
  linkcolor=uavblue,
  urlcolor=uavblue,
  citecolor=uavblue,
  pdftitle={WISP-RT Operation and Delivery Manual},
  pdfauthor={UAV-RT Team, UBC Okanagan}
}

% -------- Section styling --------
\titleformat{\section}
  {\color{uavblue}\Large\bfseries}{\thesection}{0.6em}{}
\titleformat{\subsection}
  {\color{uavblue}\large\bfseries}{\thesubsection}{0.5em}{}
\titleformat{\subsubsection}
  {\color{uavgrey}\normalsize\bfseries}{\thesubsubsection}{0.4em}{}

% -------- Header / footer --------
\pagestyle{fancy}
\fancyhf{}
\fancyhead[L]{\small\color{uavgrey}WISP-RT Operation \& Delivery Manual}
\fancyhead[R]{\small\color{uavgrey}City of Kelowna}
\fancyfoot[C]{\small\color{uavgrey}Page \thepage\ of \pageref{LastPage}}
\renewcommand{\headrulewidth}{0.4pt}
\usepackage{lastpage}

% -------- Custom boxes --------
\newtcolorbox{infobox}[1][]{
  enhanced, breakable,
  colback=uavlight, colframe=uavblue,
  boxrule=0.6pt, arc=2pt,
  left=8pt, right=8pt, top=6pt, bottom=6pt,
  title={\bfseries #1}, coltitle=white,
  colbacktitle=uavblue, sharp corners=downhill
}

\newtcolorbox{warnbox}[1][Warning]{
  enhanced, breakable,
  colback=uavlight, colframe=uavwarn,
  boxrule=0.8pt, arc=2pt,
  left=8pt, right=8pt, top=6pt, bottom=6pt,
  title={\bfseries #1}, coltitle=white,
  colbacktitle=uavwarn, sharp corners=downhill
}

\newtcolorbox{tipbox}[1][Tip]{
  enhanced, breakable,
  colback=uavlight, colframe=uavgood,
  boxrule=0.6pt, arc=2pt,
  left=8pt, right=8pt, top=6pt, bottom=6pt,
  title={\bfseries #1}, coltitle=white,
  colbacktitle=uavgood, sharp corners=downhill
}

% Photo / screenshot placeholder
% Replace with \includegraphics{filename} once photos are added.
\newcommand{\photoslot}[1]{%
  \begin{center}
  \fbox{\begin{minipage}[c][3.2cm][c]{0.85\textwidth}
    \centering\color{uavgrey}\small\itshape[PHOTO / SCREENSHOT: #1]
  \end{minipage}}
  \end{center}
}

% Checklist item
\newcommand{\chk}{\item[$\Box$]}

% -------- Title page --------
\begin{document}

\begin{titlepage}
  \centering
  \vspace*{2.2cm}
  {\color{uavblue}\rule{\textwidth}{1.2pt}}\\[0.6cm]
  {\Huge\bfseries\color{uavblue} WISP-RT}\\[0.3cm]
  {\LARGE\bfseries Operation \& Delivery Manual}\\[0.5cm]
  {\color{uavblue}\rule{\textwidth}{0.6pt}}\\[1.6cm]

  {\large Wildfire Incident Surveillance Platform, Rail-Trail}\\[0.3cm]
  {\large Autonomous Drone-in-a-Box System for Wildfire and Incident Monitoring}\\[2.5cm]

  \begin{tabular}{rl}
    \textbf{Delivered to:} & City of Kelowna \\[0.15cm]
    \textbf{Prepared by:}  & UAV-RT Team, UBC Okanagan \\[0.15cm]
    \textbf{Team:}         & \begin{tabular}[t]{@{}l@{}}
                                 Kangnan Wang \\
                                 Daniil Kovalevskiy \\
                               \end{tabular} \\[0.15cm]
    \textbf{Supervisor:}   & Dr.\ Zheng Liu \\[0.15cm]
    \textbf{Document version:} & 1.0 \\[0.15cm]
    \textbf{Date:} & September 2026 \\
  \end{tabular}

  \vfill
  {\small\color{uavgrey}This is a non-technical manual.
  The document explains how to physically set up, start, operate, monitor, and shut down
  the delivered WISP-RT system. It does not cover source code, AI algorithm
  design, or development-only procedures.}
\end{titlepage}

% -------- Table of contents --------
\tableofcontents
\newpage

% =============================================================
\section{System Overview}
% =============================================================

The WISP-RT system is an \textbf{autonomous drone monitoring station} designed
to detect and report hazards - such as fires, unattended vehicles, and
suspicious objects - across a designated area in and around the City of
Kelowna. Once installed, the system launches missions from a DJI Dock~3,
streams live video (regular and thermal) to an operator computer, runs
on-computer AI on the video, and pushes alerts to the operator's mobile
phone.

\begin{infobox}[What this system is used for]
\begin{itemize}[leftmargin=1.3em]
  \item Scheduled or on-demand aerial patrols from a fixed dock location.
  \item Live monitoring of regular and thermal video from the drone.
  \item Automatic detection of hazards (fire, smoke, people, vehicles, hot spots).
  \item Push notifications to the operator's mobile phone when an event occurs.
  \item Recording and later review of flight video and incident evidence.
\end{itemize}
\end{infobox}

\subsection{Main Hardware Components}

\begin{tabular}{@{}p{4.5cm}p{10.5cm}@{}}
\toprule
\textbf{Component} & \textbf{Purpose} \\
\midrule
DJI Matrice 4TD drone      & Aerial platform with RGB and thermal cameras. \\
DJI Dock~3                 & Weatherproof landing dock. Charges the drone, launches missions, and provides the network uplink. \\
DJI RC Plus 2 remote       & Manual override controller (used only in emergencies or for local checks). \\
Operator computer          & Windows~10/11 desktop that runs the WISP-RT operational software. \\
Uninterruptible Power Supply (UPS) & Keeps critical systems running during short-term power outages. \\
Operator mobile phone      & Receives push alerts through the \emph{ntfy} app. \\
\bottomrule
\end{tabular}

\subsection{Main Software Components}
\label{sec:overview-software}

\begin{tabular}{@{}p{4.5cm}p{10.5cm}@{}}
\toprule
\textbf{Software} & \textbf{Purpose} \\
\midrule
DJI FlightHub~2      & Cloud service where you create missions, launch the drone, and see flight records. \\
WISP-RT operational software & Web application installed on the operator computer. Shows live video, AI detections, thermal data, and the incident list. \\
ntfy mobile app      & Free notification app used to receive alerts on the operator's phone. \\
Web browser (Chrome or Edge) & Used to open both FlightHub~2 and the WISP-RT operational page. \\
\bottomrule
\end{tabular}

\subsection{Simple System Architecture}

\begin{center}
\begin{tcolorbox}[colback=uavlight, colframe=uavblue, boxrule=0.6pt, arc=2pt,
  width=0.95\textwidth]
\ttfamily\small
\begin{tabbing}
\hspace{0.5cm}\=\hspace{4.5cm}\=\hspace{4.5cm}\= \kill
\> [DJI Drone M4TD] \\
\> \ \ $\uparrow$/$\downarrow$ (radio / cellular) \\
\> [DJI Dock 3] $\longrightarrow$ (Ethernet / Cellular) $\longrightarrow$ [FlightHub 2 Cloud] \\
\> \ \ $\downarrow$ (video + telemetry) \\
\> [Local Router] \\
\> \ \ $\downarrow$ \\
\> [Operator Computer running WISP-RT] $\longrightarrow$ [ntfy] $\longrightarrow$ [Phone] \\
\end{tabbing}
\end{tcolorbox}
\end{center}

\photoslot{Overall connection diagram - photo or drawn diagram of the physical
setup showing Dock, drone, router, operator PC, and phone.}

\subsection{What Each Component is Responsible For}

\begin{description}[leftmargin=2cm, style=nextline]
  \item[Drone] Flies the mission, records video, and streams live video back through the Dock.
  \item[Dock~3] Keeps the drone charged, opens/closes the cover, provides GPS reference, and forwards data to FlightHub~2 and the operator computer.
  \item[FlightHub~2] Cloud service that stores missions, controls the Dock, and archives flight recordings.
  \item[Operator Computer] Runs the system software. Displays live video, runs the AI, and stores incident snapshots and clips.
  \item[Mobile Phone (ntfy)] Receives push notifications so the operator does not need to watch the screen at all times.
\end{description}

\newpage
% =============================================================
\section{Regulatory and Compliance Responsibility}
\label{sec:regulatory}
% =============================================================

\begin{warnbox}[The City and the operator are responsible for compliance]
Operating this system involves flying a remotely piloted aircraft in
Canadian airspace. The City of Kelowna and the pilot in command are
responsible for identifying and following the laws, regulations,
authorizations, and City policies that apply to the operation. This
manual explains how to use the delivered system; it does not restate
regulatory rules and is not a substitute for them.
\end{warnbox}

Typical authorities the City should confirm before operating include
Transport Canada (pilot certification, aircraft registration, and RPAS
operating rules), NAV CANADA (airspace coordination), and the City of
Kelowna's own bylaws and privacy policies. The City is responsible for
maintaining the compliance and insurance documents that its regulator
and internal policy require.

\begin{infobox}[Fire incident response]
Any confirmed fire detected by this system must be reported through the
City of Kelowna's existing fire-response procedure. This system is a
monitoring aid, not a replacement for professional response.
\end{infobox}

\newpage


% =============================================================
\section{City Responsibilities and Provided Services}
% =============================================================

The delivered system depends on a small number of accounts and
infrastructure items owned by the City of Kelowna. Only items confirmed
as part of the delivered configuration are listed below.

\subsection{Accounts to Maintain}

\begin{longtable}{@{}p{4.0cm}p{11.0cm}@{}}
\toprule
\textbf{Item} & \textbf{Purpose} \\
\midrule
DJI FlightHub~2 account & Cloud service that controls the Dock, stores missions, and archives flight records. Required for the Dock to operate. \\
\addlinespace
ntfy topic subscription & Free notification service. The operator's phone subscribes to the delivered topic name to receive alerts. \\
\addlinespace
\bottomrule
\end{longtable}

Any additional accounts, subscriptions, or warranty plans (for example,
manufacturer service coverage or LTE service for the Dock's cellular
fallback) are the City's decision and are not committed to by this
manual.

\subsection{Physical Infrastructure the City Must Provide}

\begin{longtable}{@{}p{3.4cm}p{7cm}p{4cm}@{}}
\toprule
\textbf{Item} & \textbf{Requirement} & \textbf{Maintained by} \\
\midrule
Dock site & Level outdoor location with clear sky view and no overhanging vegetation. & City facilities. \\
\addlinespace
AC power to Dock & Grounded AC circuit sized for the Dock's rated input. & City facilities / electrician. \\
\addlinespace
Network from Dock to operator & Wired connection sufficient for live video from the Dock to the operator computer. & City IT. \\
\addlinespace
Operator workstation & Indoor computer running the delivered WISP-RT web application in Chrome or Edge. & City IT. \\
\addlinespace
Operator mobile phone & Android or iPhone with the ntfy app installed and subscribed to the delivered topic. & City operations. \\
\bottomrule
\end{longtable}

\subsection{City Contact Roster}
\label{sec:contact-roster}

Fill in this table at delivery. Each role should have a named primary
and a backup.

\begin{longtable}{@{}p{4.5cm}p{5cm}p{5cm}@{}}
\toprule
\textbf{Role} & \textbf{Primary} & \textbf{Backup} \\
\midrule
Pilot in command               & \rule{4.5cm}{0.4pt} & \rule{4.5cm}{0.4pt} \\
Daily operator                 & \rule{4.5cm}{0.4pt} & \rule{4.5cm}{0.4pt} \\
IT / network                   & \rule{4.5cm}{0.4pt} & \rule{4.5cm}{0.4pt} \\
Facilities / power             & \rule{4.5cm}{0.4pt} & \rule{4.5cm}{0.4pt} \\
Incident response (fire)       & \rule{4.5cm}{0.4pt} & \rule{4.5cm}{0.4pt} \\
UAV-RT technical support (UBCO) & \rule{4.5cm}{0.4pt} & \rule{4.5cm}{0.4pt} \\
\bottomrule
\end{longtable}

% =============================================================
\section{Hardware Setup and Operation}
% =============================================================

This section covers the physical setup of the delivered system and the
day-to-day interaction the operator has with the hardware. The DJI User
Manuals shipped with each unit remain the authoritative reference for
detailed mechanical and electrical procedures.

\subsection{Environmental Operating Limits}

Missions must remain within the platform's rated environmental
envelope. These values come from the delivered platform's
specifications.

\begin{tabular}{@{}p{4.5cm}p{4.8cm}p{4.8cm}@{}}
\toprule
\textbf{Parameter} & \textbf{DJI Dock 3} & \textbf{Matrice 4TD} \\
\midrule
Ingress protection    & IP56              & IP55 \\
Operating temperature & -30 to +50\,\textdegree{}C & -20 to +50\,\textdegree{}C \\
Wind resistance       & Landing: 12~m/s   & Operation and takeoff/landing: 12~m/s \\
\bottomrule
\end{tabular}

\subsection{Site and Base Location}

The planned final base for the delivered system is the City of Kelowna
Public Works Yard at 1495 Hardy Street. The site must provide:

\begin{itemize}[leftmargin=1.5em]
  \item A level outdoor location for the Dock with a clear view of the sky.
  \item No overhanging vegetation, wires, or structures above the Dock.
  \item Grounded AC power sized for the Dock's rated input.
  \item A wired network path from the Dock to the operator workstation
        sufficient for continuous live video.
\end{itemize}

\photoslot{Photo of the Dock installed on its final pad at the Public Works
Yard, showing power and network connections.}

\subsection{Hardware Installation Steps}

The steps below are performed once, at delivery, by the UAV-RT team
together with City facilities and IT staff. They are included here so
that the City has a record of what was installed and can verify the
setup later.

\subsubsection{Step 1 -- Dock placement and power}

\begin{enumerate}[leftmargin=1.5em]
  \item Position the DJI Dock 3 on its prepared level pad at 1495 Hardy Street.
  \item Confirm the Dock is level using the built-in indicator on the top surface.
  \item Route the AC power cable to the Dock's power input; do not power on yet.
  \item Confirm no vegetation, wires, or roof overhangs sit above the Dock cover.
\end{enumerate}

\photoslot{Photo of the Dock 3 in its final position with the level indicator visible.}

\subsubsection{Step 2 -- Network cabling}

\begin{enumerate}[leftmargin=1.5em]
  \item Route the Ethernet cable from the Dock's network port into the
        operator building.
  \item Terminate the cable at the operator workstation's network switch or router.
  \item Confirm the cable is properly seated at both ends before the Dock
        is powered on.
\end{enumerate}

\photoslot{Photo of the Dock network port with the Ethernet cable connected.}

\subsubsection{Step 3 -- First power-on and Dock binding}

\begin{enumerate}[leftmargin=1.5em]
  \item Apply AC power to the Dock. The internal indicator lights turn on.
  \item Wait for the Dock to complete its self-check (typically several minutes).
  \item Bind the Dock to the City's DJI FlightHub 2 account following the
        prompts in the FlightHub 2 web interface.
  \item Confirm the Dock appears as \textbf{Online} in the FlightHub 2
        device list.
\end{enumerate}

\photoslot{Screenshot of the FlightHub 2 device list showing the Dock as Online.}

\subsubsection{Step 4 -- Aircraft placement}

\begin{enumerate}[leftmargin=1.5em]
  \item Open the Dock cover using the FlightHub 2 remote command.
  \item Place the DJI Matrice 4TD on the Dock's landing pad according to the
        alignment marks.
  \item Close the Dock cover from FlightHub 2.
  \item Confirm the drone begins charging (charging indicator in FlightHub 2).
\end{enumerate}

\photoslot{Photo of the Matrice 4TD seated on the Dock pad.}

\subsubsection{Step 5 -- Operator computer and software}

\begin{enumerate}[leftmargin=1.5em]
  \item Set up the operator computer at the workstation with monitor,
        keyboard, mouse, and network connection.
  \item The UAV-RT team installs and configures the WISP-RT operational
        software on the computer.
  \item Two desktop shortcuts are created: \textbf{WISP-RT} and \textbf{FlightHub 2}.
  \item Confirm both shortcuts open successfully in Chrome.
\end{enumerate}

\photoslot{Screenshot of the Windows desktop showing the two shortcuts.}

\subsubsection{Step 6 -- Mobile phone and ntfy}

\begin{enumerate}[leftmargin=1.5em]
  \item Install the free \textbf{ntfy} app from Google Play (Android) or
        the App Store (iPhone).
  \item Open ntfy and tap \textbf{Subscribe to topic}.
  \item Enter the topic name provided by the UAV-RT team at delivery.
  \item Tap \textbf{Subscribe}.
  \item Send a test alert from the WISP-RT Live view and confirm the phone
        receives it.
\end{enumerate}

\photoslot{Screenshot of the ntfy app subscribe screen with the delivered topic name filled in.}

\subsubsection{Step 7 -- End-to-end verification}

Before handover, the following are verified together:

\begin{itemize}[leftmargin=1.5em]
  \item Dock shows Online in FlightHub 2.
  \item Aircraft shows Ready in FlightHub 2.
  \item WISP-RT Overview page shows all System Status cells green.
  \item Live RGB and thermal video appear in the WISP-RT Live view.
  \item A test push is received on the operator's phone.
  \item A short verification flight completes normally and lands in the Dock.
\end{itemize}

\photoslot{Photo of the completed installed system - Dock on its pad, operator computer at the workstation, phone in hand.}

\subsection{Equipment the Operator Interacts With}

Day-to-day, the operator interacts only with the following items. All
other setup is complete after installation.

\begin{longtable}{@{}p{4.8cm}p{10.2cm}@{}}
\toprule
\textbf{Item} & \textbf{What the operator does with it} \\
\midrule
DJI Matrice 4TD drone       & No direct handling. Confirm it is seated in the Dock before each mission. \\
DJI Dock 3                  & No direct handling. Confirm the cover is closed between missions. \\
DJI RC Plus 2 remote controller & Used only for local overrides or supervised training. Stored powered off. \\
Operator computer           & Main daily tool. Runs WISP-RT and FlightHub~2 in a web browser. \\
Operator monitor, keyboard, mouse & Standard desktop use. \\
Operator mobile phone       & Receives ntfy alerts. Keep charged and within reach. \\
\bottomrule
\end{longtable}

\subsection{Daily System Start-Up (Operator View)}

The Dock, drone, router, and UPS are designed to run 24/7 and will
normally already be powered on when the operator arrives. The daily
start-up for the operator is limited to:

\begin{enumerate}[leftmargin=1.5em]
  \item Wake or power on the operator computer.
  \item Open the \textbf{WISP-RT} desktop shortcut (loads the
        WISP-RT Overview page).
  \item Open the \textbf{FlightHub 2} desktop shortcut and log in.
  \item Confirm the System Status card in WISP-RT shows all cells green
        and the overall pill reads ``healthy''.
  \item Confirm FlightHub~2 shows the Dock as ``Online'' and the drone as
        ``Ready''.
\end{enumerate}

\photoslot{Screenshot of the WISP-RT Overview page with all System Status
cells green and the overall pill reading ``healthy''.}

If any indicator is not green, do \textbf{not} launch a mission. Follow the
steps in Section~\ref{sec:exceptions} or contact the UAV-RT support team.

\subsection{Daily System Shut-Down (Operator View)}

Normal daily shut-down is \textbf{not required}. The system is designed to
stay online continuously. The operator only needs to:

\begin{enumerate}[leftmargin=1.5em]
  \item Confirm the drone is seated inside the Dock and the cover is closed.
  \item Confirm the WISP-RT Overview page shows no in-progress mission.
  \item Log out or lock the operator computer.
\end{enumerate}

Full power-off of the Dock and the drone is only required for maintenance
or transport, and should be coordinated with UAV-RT support.

\newpage
% =============================================================
\section{Software Preparation and Access}
% =============================================================

The four software components used day to day are introduced in
Section~\ref{sec:overview-software} of the System Overview. This section
covers installation, log-in, and the network connections between them.

\subsection{Installation and Login}

\subsubsection{Operator Computer (already installed by UAV-RT team)}

The UAV-RT team installs and configures the WISP-RT software directly on the
City's designated computer before handover. The operator does not need to
install anything on the computer. Two desktop shortcuts are provided:

\begin{itemize}[leftmargin=1.5em]
  \item \textbf{WISP-RT} - opens the WISP-RT Overview page in Chrome.
  \item \textbf{FlightHub 2} - opens DJI FlightHub~2 in Chrome.
\end{itemize}

\photoslot{Screenshot of the Windows desktop showing both shortcuts.}

To log in to FlightHub~2, click the ``FlightHub 2'' shortcut. When the DJI
sign-in page appears, enter the account credentials provided in the sealed
credentials envelope that accompanies this manual.

\begin{warnbox}[Keep credentials safe]
Do not share the FlightHub~2 credentials over email or public chat.
\end{warnbox}

\subsubsection{Mobile Phone (ntfy)}

Install ntfy on the operator's mobile phone:

\begin{enumerate}[leftmargin=1.5em]
  \item On the phone, open Google Play (Android) or App Store (iPhone).
  \item Search for \textbf{``ntfy''} (developer name: Philipp Heckel).
  \item Install the app. It is free and requires no account.
  \item Open ntfy $\rightarrow$ ``Subscribe to topic''.
  \item Enter the topic name provided by the UAV-RT team (for example, \texttt{dji\_alert}).
  \item Tap ``Subscribe''.
  \item Send a test alert from the WISP-RT Live view to confirm the phone receives it.
\end{enumerate}

\photoslot{Screenshot of the ntfy app subscribe screen with the topic name
filled in.}

\subsection{Basic Connection Setup}

The connections below are configured once by the UAV-RT team at delivery.
The operator does not usually need to change them.

\begin{tabular}{@{}p{5cm}p{10cm}@{}}
\toprule
\textbf{Connection} & \textbf{Configured value / requirement} \\
\midrule
Dock $\leftrightarrow$ FlightHub~2 & Dock is bound to the City's FlightHub~2 project. \\
FlightHub 2 $\leftrightarrow$ WISP-RT    & Media and telemetry sharing. \\
WISP-RT $\leftrightarrow$ Computer    & Local web app. Open via the \textbf{WISP-RT} desktop shortcut. \\
Phone $\leftrightarrow$ ntfy       & Subscribed to the delivered topic name. \\
\bottomrule
\end{tabular}

\newpage
% =============================================================
\section{Pre-Operation Check}
\label{sec:preop}
% =============================================================

Complete this check \textbf{before every mission}. A printable one-page
checklist is provided in Section~\ref{sec:checklist}.

\subsection{Device and Connection Check}

Open the WISP-RT Overview page. Confirm each of the following:

\begin{itemize}[leftmargin=1.5em]
  \chk Dock~3 is Online (green dot in System Status).
  \chk Drone is Online (green dot in Aircraft telemetry card).
  \chk Camera (M4TD) is Online.
  \chk Live streams counter reads at least the minimum for the deployed cameras (thermal and dock feeds ready).
  \chk All System Status cells on the WISP-RT Overview page are green.
  \chk FlightHub~2 dashboard shows the Dock as ``Online''.
\end{itemize}

\photoslot{Screenshot of the WISP-RT System Status card with all cells
green.}

\subsection{Function Check}

Go to the WISP-RT ``Live view'' page and confirm each of the following:

\begin{itemize}[leftmargin=1.5em]
  \chk Regular (RGB) live video appears within 10~seconds.
  \chk Thermal live video appears (switch to thermal tile).
  \chk Temperature reading shows a value (top-left of thermal tile).
  \chk The AI worker card shows ``Running'' with a green dot.
  \chk Press the ``Send test push'' button - phone receives an ntfy alert.
\end{itemize}

\subsection{Ready-to-Operate Check}

Complete every item below before launch. The regulatory items below
must be satisfied per the applicable rules and the City's own
procedures - see Section~\ref{sec:regulatory}.

\begin{itemize}[leftmargin=1.5em]
  \chk Drone battery at least 60~\%.
  \chk Weather within DJI operating limits (no precipitation, icing, or wind above the M4TD manufacturer limit).
  \chk NOTAMs reviewed for the operating area.
  \chk Any required airspace authorization is in place and accessible to the pilot.
  \chk Site assessment complete; people-control measures in place.
  \chk Pilot certification, recency evidence, and aircraft registration accessible to the pilot.
  \chk Mission plan selected in FlightHub~2.
  \chk Lost-link / RTH behaviour reviewed for this specific site.
  \chk Emergency contact list visible at the operator station.
\end{itemize}

\newpage
% =============================================================
\section{Mission Operation}
% =============================================================

\subsection{Create a Flight Route}

\begin{enumerate}[leftmargin=1.5em]
  \item Open FlightHub~2 using the desktop shortcut.
  \item In the left menu, click \textbf{Flight Route Library}.
  \item Click \textbf{Create Route}.
  \item Choose the route type: \textbf{Waypoint}, \textbf{Mapping}, or \textbf{Oblique}.
        For general patrols, use \textbf{Waypoint}.
  \item On the map, click each waypoint the drone should visit.
  \item For each waypoint, set the parameters required by the mission:
    \begin{itemize}[leftmargin=1.2em]
      \item Altitude above ground.
      \item Speed.
      \item Actions (hover, photo, gimbal angle, etc.).
    \end{itemize}
        Use the altitude and speed values established during commissioning
        for the selected route; they are not fixed defaults in this manual.
  \item Click \textbf{Save} and give the route a clear name, e.g.
        ``\texttt{Downtown\_Patrol\_A}''.
\end{enumerate}

% \photoslot{Screenshot of the FlightHub~2 route creation view with waypoints
% placed on the map.}

\begin{figure}[htp]
    \centering
    \includegraphics[width=16cm]{images/Screenshot 2026-09-12 at 11.19.39.png}
    \caption{FlightHub 2 Flight Route Creation}
\end{figure}

\subsection{Start and Execute a Mission}

\begin{enumerate}[leftmargin=1.5em]
  \item Open FlightHub~2 $\rightarrow$ \textbf{Projects} $\rightarrow$ select your project.
  \item Click \textbf{Task Plan Library} $\rightarrow$ \textbf{Create Task}.
  \item Select the route saved in the previous step.
  \item Review the pre-flight checklist that FlightHub~2 shows automatically
        (battery, weather, RTH altitude).
  \item Click \textbf{OK}.
  \item Watch the Dock cover open, the drone lift off, and the mission counter start.
  \item Switch to the WISP-RT ``Live view'' page to watch the flight in real time.
\end{enumerate}

% \photoslot{Screenshot of the FlightHub~2 ``Execute Task'' confirmation dialog.}

\begin{figure}[htp]
    \centering
    \includegraphics[width=16cm]{images/Screenshot 2026-09-12 at 11.31.16.png}
    \caption{FlightHub 2 Mission Execution}
\end{figure}

\subsection{Pause, Stop, and Complete a Mission}

\subsubsection{Pause}

In FlightHub~2, click the yellow \textbf{Pause} button on the mission control
strip. The drone will hover in place.

\subsubsection{Manual Stop}

Click the red \textbf{Stop} button. The drone will follow its programmed
Return-to-Home procedure and land in the Dock.

\subsubsection{Emergency Return}

Click \textbf{Return to Home (RTH)}. This is the same as Stop but forces
the drone to climb to the RTH altitude before returning.

\subsubsection{Mission Complete}

When the drone lands and the Dock cover closes, confirm:

\begin{itemize}[leftmargin=1.5em]
  \item FlightHub~2 shows the task as ``Complete''.
  \item The WISP-RT Overview page shows one more ``Flight today'' count.
  \item Any incidents flagged during the flight appear in the Incidents list.
\end{itemize}

\newpage
% =============================================================
\section{Monitoring, Alerts, and Records}
% =============================================================

\subsection{Live Monitoring}

Open the WISP-RT ``Live view'' page. You will see:

\begin{itemize}[leftmargin=1.5em]
  \item \textbf{Primary video pane} - large window with overlays (AI boxes, thermal palette, temperature reading).
  \item \textbf{Camera tiles} on the right - click any tile to switch the primary view.
  \item \textbf{Flight status bar} - altitude, speed, battery, wind, and heading.
\end{itemize}

% \photoslot{Annotated screenshot of the Live view showing the primary pane,
% camera tiles, split-view button, and telemetry strip.}

\begin{figure}[htp]
    \centering
    \includegraphics[width=16cm]{images/Screenshot 2026-09-12 at 11.48.46.png}
    \caption{WISP-RT Live View}
\end{figure}

\subsection{AI and Thermal Information}

\subsubsection{Where AI Results Appear}

\begin{itemize}[leftmargin=1.5em]
  \item Coloured boxes on the primary video pane - one box per detected object.
  \item Header \textbf{bell icon} - unread alert count.
  \item Overview page \textbf{Incidents list} - all persisted alerts of the day, each with a snapshot photo and a 5-second video clip.
\end{itemize}

\subsubsection{How the System Decides When to Alert}

The delivered system combines object detection with thermal information. The
logic follows the day/night sensing priority of the Matrice 4TD:

\begin{itemize}[leftmargin=1.5em]
  \item \textbf{Daytime} -- detection is based mainly on the RGB (visible-light) camera.
  \item \textbf{Nighttime} -- detection is based mainly on the IR (thermal) camera.
  \item Detected objects are considered together with the temperature information
        from the same scene.
  \item An \textbf{alert is generated} when either:
    \begin{itemize}[leftmargin=1.2em]
      \item the object's temperature is abnormally above its normal range for
            that class of object, or
      \item fire or smoke is directly detected.
    \end{itemize}
  \item Otherwise, the observation is treated as normal.
\end{itemize}

\begin{infobox}[An object label alone is not an incident]
A bounding box on a person, vehicle, or tent does not by itself indicate a
fire incident. The alert only fires when the thermal reading is out of the
normal range for the observed target, or when fire or smoke is detected
directly. The operator still reviews the live evidence before deciding on a
response.
\end{infobox}

\subsubsection{Detection Classes}

The deployed detector recognizes ten classes:

\begin{tabular}{@{}p{3.5cm}p{11.5cm}@{}}
\toprule
\textbf{Class} & \textbf{Meaning for the operator} \\
\midrule
Fire            & Visible flame or extreme heat source. \textbf{Immediate response required.} \\
Smoke           & Smoke plume detected. Confirm and log. \\
Person          & Person present in the monitored area. \\
Vehicle         & Car, truck, or similar in the monitored area. \\
Tree            & Vegetation context. \\
Dry\ grass      & Combustible vegetation. Alert only if the temperature is abnormal for the class. \\
Tent / Tarp     & Temporary shelter. Alert only if the temperature is abnormal for the class. \\
Plastic waste   & Debris. Alert only if the temperature is abnormal for the class. \\
Metal fixture   & Fixed infrastructure. Context class. \\
Green grass     & Non-combustible vegetation. Context class. \\
\bottomrule
\end{tabular}

\subsection{Mobile Alerts (ntfy)}

When an alert fires:

\begin{enumerate}[leftmargin=1.5em]
  \item The phone shows a notification with:
    \begin{itemize}[leftmargin=1.2em]
      \item Class (fire, smoke, person, etc.).
      \item Temperature (if applicable).
      \item Short description (``$\times$3 in 10\,s'' if several detections were bundled).
    \end{itemize}
  \item Tapping the notification opens the WISP-RT Live view page.
  \item The operator should then:
    \begin{itemize}[leftmargin=1.2em]
      \item Open the WISP-RT Overview page.
      \item Find the alert in the Incidents list.
      \item Click the thumbnail to view the snapshot and 5-second clip.
      \item Decide on next action (log, forward, dispatch, ignore).
    \end{itemize}
\end{enumerate}

% \photoslot{Photo of a phone showing a fire alert notification and screenshot
% of the WISP-RT Incidents list with an unread alert highlighted.}

\subsection{Flight Records and Video}

\subsubsection{In FlightHub~2}

\begin{enumerate}[leftmargin=1.5em]
  \item Open FlightHub~2 $\rightarrow$ \textbf{Missions} $\rightarrow$ \textbf{History}.
  \item Filter by date or Dock.
  \item Click a mission to see:
    \begin{itemize}[leftmargin=1.2em]
      \item Flight path replay on the map.
      \item Recorded video (RGB and thermal).
      \item Telemetry graphs.
    \end{itemize}
  \item Click \textbf{Download} to save video locally if needed.
\end{enumerate}

\subsubsection{In System}

\begin{itemize}[leftmargin=1.5em]
  \item \textbf{Overview page} $\rightarrow$ \textbf{Incidents list}: every persisted alert with photo and 5~s clip.
  \item \textbf{Footage page}: thumbnail browser of all captured incident media, grouped by date.
  \item Click any card to open the preview modal (photo + inline video).
  \item Use the \textbf{Forward} button to send selected incidents to the phone via ntfy.
  \item Use the \textbf{Clear} button to remove old incidents from the list.
\end{itemize}

% \photoslot{Screenshot of the WISP-RT Incidents list with two items checked and
% the ``Forward'' button highlighted.}

\begin{figure}[htp]
    \centering
    \includegraphics[width=16cm]{images/Screenshot 2026-09-12 at 12.22.59.png}
    \caption{WISP-RT Incident Handling}
\end{figure}

\newpage
% =============================================================
\section{Routine Maintenance and Firmware Updates}
\label{sec:maintenance}
% =============================================================

This section covers the routine checks the daily operator performs and
the firmware-update workflow in FlightHub~2. Manufacturer-defined
service intervals for the drone, batteries, and Dock consumables are
specified in the DJI User Manuals and take precedence over any interval
suggested here.

\subsection{Operator Daily Checks}

\begin{itemize}[leftmargin=1.5em]
  \item Visual walk-around of the Dock: cover closed, no debris, no water pooling on the pad.
  \item WISP-RT Overview page shows all System Status cells green.
  \item Confirm ntfy alerts are received on the operator phone using the test-push button.
\end{itemize}

\subsection{Operator Weekly Checks}

\begin{itemize}[leftmargin=1.5em]
  \item Inspect drone propellers for chips, cracks, or bent tips. Replace
        any propeller with visible damage.
  \item Wipe the RGB and thermal camera lenses with the DJI cleaning cloth.
  \item Check FlightHub~2 for firmware update notifications and, if one is
        available, coordinate the update as described below.
\end{itemize}

\photoslot{Photo of a healthy propeller next to one with a small tip chip
that requires replacement.}

\subsection{DJI-Specified Service}

Battery cycle count and health, Dock desiccant replacement, cabling
inspection, and manufacturer service intervals are defined by DJI. The
City's designated technician follows the schedules in the
\emph{DJI Matrice 4TD User Manual} and \emph{DJI Dock 3 User Manual};
they are not restated here to avoid conflicting with the current
manufacturer guidance.

\subsection{Firmware Updates in FlightHub~2}

DJI releases firmware updates for the drone, Dock, and remote controller
several times a year. Updates fix bugs, patch security issues, and
occasionally add capabilities. The City's technician should check for
updates \textbf{weekly} and apply them promptly.

\subsubsection{Weekly Firmware Check}

\begin{enumerate}[leftmargin=1.5em]
  \item Open FlightHub~2 in the operator computer's browser and sign in.
  \item Go to \textbf{Devices} $\rightarrow$ select the Dock.
  \item Look at the device summary panel. If a firmware update is
        available, a coloured \textbf{Update} banner or badge appears next to the device.
  \item Click the banner to see the release notes and target version.
\end{enumerate}

% \photoslot{Screenshot of the FlightHub~2 Devices panel showing an
% ``Update available'' badge next to the Dock.}

\begin{figure}[htp]
    \centering
    \includegraphics[width=16cm]{images/Screenshot 2026-09-12 at 12.52.50.png}
    \caption{FlightHub 2 Firmware Updates}
\end{figure}

\subsubsection{Applying a Firmware Update}

\begin{warnbox}[Before you start]
\begin{itemize}[leftmargin=1.3em]
  \item Confirm no mission is in progress or scheduled to start within the next 60~minutes.
  \item Confirm the drone is seated in the Dock and the cover is closed.
  \item Confirm the drone battery is above 50~\%.
  \item Confirm the ground internet connection is stable (Ethernet, not just LTE fallback).
  \item Never cut power to the Dock or drone during an update.
\end{itemize}
\end{warnbox}

\begin{enumerate}[leftmargin=1.5em]
  \item In FlightHub~2, open the Dock's device page and click \textbf{Firmware Update} $\rightarrow$ \textbf{Update Now}.
  \item FlightHub~2 downloads the update to the Dock (typically 5--20~minutes depending on size).
  \item The Dock then flashes the drone and the Dock modules. The drone and Dock will reboot at least once.
  \item Wait until FlightHub~2 shows the new version as ``Installed'' and the device status returns to \textbf{Online}.
  \item Open the WISP-RT Overview page and confirm the System Status card is fully green again. Some cells may take a few minutes to reconnect after the reboot.
  \item Fly a short test mission (or launch and land at a low altitude) to confirm behaviour is unchanged.
\end{enumerate}

\subsubsection{If a Firmware Update Fails}

\begin{itemize}[leftmargin=1.5em]
  \item Do \textbf{not} power-cycle the Dock or drone. FlightHub~2 usually retries automatically.
  \item If FlightHub~2 reports the device as ``offline for update'' for more than 30~minutes, contact the UAV-RT support team before doing anything physical.
  \item Never attempt to downgrade firmware unless DJI or UAV-RT support explicitly instructs you to.
\end{itemize}

\subsection{WISP-RT Software Updates}

The WISP-RT operational software is maintained by the UAV-RT team. New
versions are delivered as a scheduled visit or a remote-assistance
session. The operator does not need to install or patch WISP-RT.

\newpage
% =============================================================
\section{Basic Exception and Safety Handling}
\label{sec:exceptions}
% =============================================================

\subsection{Connection Loss / Return-to-Home}

If the drone loses communication with the Dock during flight:

\begin{itemize}[leftmargin=1.5em]
  \item The drone automatically enters DJI Return-to-Home behaviour.
  \item It will climb to the pre-set RTH altitude, fly back to the Dock, and land.
  \item The operator should:
    \begin{enumerate}[leftmargin=1.2em]
      \item Watch the FlightHub~2 map until the drone reaches the Dock.
      \item Not attempt to re-launch until the connection is restored.
    \end{enumerate}
\end{itemize}

\subsection{Mission Interruption}

If a mission is interrupted (weather, obstacle, operator decision):

\begin{enumerate}[leftmargin=1.5em]
  \item Click \textbf{Pause} in FlightHub~2 if the drone should hover.
  \item Click \textbf{RTH} or \textbf{Stop} if the drone must come home.
  \item Do not power the Dock off while the drone is in the air.
  \item Once the drone has landed, review the FlightHub~2 log for the reason.
\end{enumerate}

\subsection{Software / AI Unavailable}

If the WISP-RT Overview page shows red or ``offline'' dots for AI worker,
Live streams, or Thermal SDK:

\begin{itemize}[leftmargin=1.5em]
  \item Do \textbf{not} rely on the affected function during the flight.
  \item Continue the mission only if visual monitoring in FlightHub~2 alone is acceptable.
  \item If AI monitoring is critical, pause or postpone the mission.
  \item Contact the UAV-RT technical support team (contact details on the
        inside cover of this manual).
\end{itemize}

\begin{tipbox}[Quick recovery]
For most WISP-RT issues, refreshing the browser page (F5) or restarting the
computer resolves the problem. Do not re-install any software - contact
support instead.
\end{tipbox}

\subsection{Alert Response}

After receiving a warning on the phone:

\begin{enumerate}[leftmargin=1.5em]
  \item Open the WISP-RT Live view (tap the notification, or click the desktop shortcut).
  \item Confirm the event on live video.
  \item Open the Overview page and find the incident.
  \item Review the snapshot and 5~s clip.
  \item Decide whether to:
    \begin{itemize}[leftmargin=1.2em]
      \item Forward to emergency services or the City response team.
      \item Save the record for follow-up.
      \item Dismiss (mark as reviewed).
    \end{itemize}
  \item Follow the City's internal incident response procedure.
\end{enumerate}

\begin{warnbox}[When in doubt, escalate]
For any suspected fire alert, follow the City's fire response procedure
first. This system is a monitoring aid, not a replacement for professional
response.
\end{warnbox}

\newpage
% =============================================================
\section{Quick Operation Guide}
\label{sec:checklist}
% =============================================================

The one-page guide below covers a normal operating day. Print this page and
keep it near the operator computer.

\begin{tcolorbox}[colback=uavlight, colframe=uavblue, boxrule=0.8pt, arc=3pt,
  title={\bfseries Daily WISP-RT Operating Workflow}, coltitle=white,
  colbacktitle=uavblue, fonttitle=\Large]

\textbf{1. Confirm hardware in place}
\begin{itemize}[leftmargin=1.5em, itemsep=0.1em]
  \item Dock, drone, operator computer, router, and phone all present and undamaged.
  \item UPS ``online'' indicator is green; router internet indicator is on.
\end{itemize}

\textbf{2. Start the operator computer}
\begin{itemize}[leftmargin=1.5em, itemsep=0.1em]
  \item The Dock, drone, router, and UPS run 24/7 - do not power-cycle them.
  \item Wake or log in the operator computer; open the \textbf{WISP-RT} and \textbf{FlightHub 2} desktop shortcuts.
\end{itemize}

\textbf{3. Check system status}
\begin{itemize}[leftmargin=1.5em, itemsep=0.1em]
  \item WISP-RT Overview page: all System Status cells green.
  \item FlightHub~2: Dock and drone show ``Online''.
  \item Phone: ntfy subscribed to the correct topic.
\end{itemize}

\textbf{4. Create or select mission}
\begin{itemize}[leftmargin=1.5em, itemsep=0.1em]
  \item FlightHub~2 $\rightarrow$ Flight Route Library $\rightarrow$ pick or create a route.
\end{itemize}

\textbf{5. Start mission}
\begin{itemize}[leftmargin=1.5em, itemsep=0.1em]
  \item FlightHub~2 $\rightarrow$ Create Task $\rightarrow$ Immediate Task $\rightarrow$ Execute.
\end{itemize}

\textbf{6. Monitor flight and AI results}
\begin{itemize}[leftmargin=1.5em, itemsep=0.1em]
  \item WISP-RT Live view: watch video, AI boxes, and temperature overlay.
  \item Overview page: watch the Incidents list update.
\end{itemize}

\textbf{7. Receive and check alerts}
\begin{itemize}[leftmargin=1.5em, itemsep=0.1em]
  \item Phone ntfy notification $\rightarrow$ open WISP-RT $\rightarrow$ confirm on Live view $\rightarrow$ review Incidents list.
\end{itemize}

\textbf{8. Complete mission}
\begin{itemize}[leftmargin=1.5em, itemsep=0.1em]
  \item Drone lands automatically. FlightHub~2 shows ``Complete''.
\end{itemize}

\textbf{9. Review records}
\begin{itemize}[leftmargin=1.5em, itemsep=0.1em]
  \item FlightHub~2 $\rightarrow$ Missions $\rightarrow$ History.
  \item WISP-RT Footage page for incident snapshots and clips.
\end{itemize}

\textbf{10. End of day}
\begin{itemize}[leftmargin=1.5em, itemsep=0.1em]
  \item Confirm drone is seated with cover closed. Confirm no mission in progress. Log out or lock the computer.
  \item Full system shutdown is \emph{not} required day-to-day. Contact UAV-RT support if maintenance or transport requires powering the Dock off.
\end{itemize}

\end{tcolorbox}

\vspace{0.6cm}

\begin{tcolorbox}[colback=uavlight, colframe=uavwarn, boxrule=0.6pt,
  title={\bfseries Pre-Mission Checklist}, coltitle=white,
  colbacktitle=uavwarn]
\begin{itemize}[leftmargin=1.5em, itemsep=0.1em]
  \chk Dock, drone, and camera all Online.
  \chk Battery $\geq$ 60~\%.
  \chk Live RGB video visible.
  \chk Live thermal video visible.
  \chk AI worker Running.
  \chk Test push received on phone.
  \chk Weather within safe limits.
  \chk Mission selected in FlightHub~2.
\end{itemize}
\end{tcolorbox}

\vspace{0.6cm}

\begin{center}
\small\color{uavgrey}\itshape
End of Operation \& Delivery Manual. \\
For technical support, contact the UAV-RT team at UBC Okanagan.
\end{center}

\end{document}